\section*{The macroscopic interval likelihood and its approximations}
% Owner: papers/_program/nomenclature.md (letters, lattice, ladder, region names) and
% papers/_program/program.md §1 (body roster). Presents the root question (are the gating
% fluctuations modelled at all?) with least squares at the bottom of the cost ladder, and the
% lattice of likelihoods that hangs from it, with the endpoint ladder as the presentation.
% Do NOT restate: model scope (program.md §7), the novelty claim, the diagnostic definitions
% (03_diagnostics.tex), or run values and the noise parameter (06_methods.tex). The boundary
% moments are given here in the form the code computes them (spectral, E_2 and E_3); what stays for
% the Supplement is their derivation and the coincidence branches at index level.
% AMENDED 2026-08-06: this line used to send E_2 and E_3 to Methods and the Supplement, and neither
% carried them, so the two kernels were nowhere. They are in the body now, at matrix level.
% COMPRESSED 2026-08-09 (length pass). 606 words -> 338. What went: the two paragraphs of the box
% that COMMENTED on the convention instead of stating it ("Letters name states and the arrow names
% the time relation..." and "The arrow says what the object is, not how it was obtained..."); the
% dual-reading rule for pair-indexed quantities, with it the body's definitions of the entrywise
% product and the ones vector, since no K^2 object is formed in the body any more and both symbols
% are defined where they are used (08_appendix_derivation.tex, at the spectral kernels and at
% Eq. tilde, and 06_methods.tex before the dictionary table); the binomial count of occupancy
% vectors; and the closing "The rest of this section unpacks that sentence". Kept whole: the two-
% state scope, both obstructions and their asymmetry (cost vs no closed form), and the availability
% of the low-order moments. Added: "of duration $\Delta$" and "from state $i$ to state $j$", four
% words each, because the cut paragraph was where the reader learned that the letters name states,
% and $\Delta$ is not defined until the next subsection.

A macroscopic current is the summed signal of a population of ion channels, each an
independent continuous-time Markov chain over a few conformational states. $N_{\mathrm{ch}}$
is the number of channels and $K$ the number of kinetic states of one channel ($K=2$ % src: papers/_program/program.md §7 (model scope: two states)
throughout this paper, carried symbolically). Each channel has a generator (rate) matrix
$\mathbf{Q}\in\mathbb{R}^{K\times K}$, hence the row-stochastic transition matrix
$\mathbf{P}(t)=e^{\mathbf{Q}t}$ over an elapsed time $t$, with entries $P_{i\to j}(t)$ from
state $i$ to state $j$, and a conductance $\gamma_k$ in state $k$, collected in
$\bm{\gamma}\in\mathbb{R}^{K}$.

\begin{tcolorbox}[title={Three marks, and each one says one thing}, colback=gray!5, colframe=gray!55]
\textbf{An arrow means the object spans the acquisition window}, of duration $\Delta$:
$P_{i\to j}(\Delta)$ is a transition across it, $\bar\gamma_{i\to j}$ an average over it,
$\bar y_{0\to\Delta}$ a recorded value formed from it. \textbf{No arrow means one instant}:
$\bm{\mu}_0,\bm{\Sigma}_0$ are the occupancy mean and covariance where the window opens,
$\bm{\mu}_\Delta,\bm{\Sigma}_\Delta$ where it closes. \textbf{A dot marks a collapsed end}:
$a_{i\to\cdot}:=\sum_j P_{i\to j}(\Delta)\,a_{i\to j}=\mathbb{E}[a\mid X_0=i]$, as in
$\bar\gamma_{i\to\cdot}$; for a probability the weight is already inside and the collapse is the
plain margin. Across a recording a subscript $t$ indexes the window, as in $\bar y_t$; inside a window the two
ends are $0$ and $\Delta$, and a collapsed object is defined the moment its pair-resolved version is.
\end{tcolorbox}
% NOTATION, aligned 2026-08-09 with the published form. The supplementary information of
% moffatt2025bayesian, section S1.3, writes P_{i->j}(t), gamma_{i->j}(t), gamma_0 = (E[gamma_{i->.}])_i,
% y^pred_{0->t}, y^obs_{0->t} and K^mu_{0->t}: lowercase, i and j, the dot for the collapse, and the
% arrow on every spanning object INCLUDING the gain. Thirteen pages with no uppercase Gamma, no
% i_0/i_t, no "bnd" and no "prop". This section had departed from that on all four counts; the pass
% of 2026-08-09 returns it. VERIFIED first-hand against
% docs/bibliography/Moffatt_PierdominiciSottile_2025_Functional_Asymmetry_P2X2_CommBiol_SI.pdf,
% Eqs. S1-S6.
% THREE DELIBERATE DEPARTURES, kept:
%   (i) the bar on gamma-bar. The SI writes gamma_{i->j} bare because nothing competes with it
%       there; this section uses bold gamma for the instantaneous conductance in the same
%       equations, so the bar is a local disambiguation and not a departure of substance.
%  (ii) the gain is g and not K^mu, and it is not called a Kalman gain. De-Kalmaned 2026-08-06
%       with the reason recorded below Eq. sig-post.
% (iii) the SI's E_i for the collapsed second moment is not used: E_2 and E_3 are the divided
%       differences here, and a third E would collide.

Writing the likelihood of a recorded sample exactly, given the kinetic parameters and the
data so far, calls for two separate objects. The first is a microscopic description that
keeps each channel discrete, the ensemble an integer occupancy vector; it is explicit, and
the recursion over it is the microscopic filter of \citet{moffatt2007estimation}, but each
sample costs the square of the number of occupancy vectors and is out of reach at
macroscopic channel counts. The second is the law of the conductance averaged over the
acquisition window given that description, which has no closed expression at any channel
count, although its low-order moments do. The two obstructions differ in kind, cost for the
first and no closed form at all for the second. Every likelihood that runs at these channel
counts replaces both objects with Gaussians, and the members of the family studied here differ in which
object the second Gaussian describes, what each is conditioned on and how its variance is
accounted.
% CORRECTED 2026-08-06. The paragraph read "Its exact form would require the distribution of
% the interval-averaged conductance of the whole population, which has no closed expression",
% which names one of the two obstructions and attributes the whole intractability to it. The
% exact likelihood needs the microscopic (occupancy-vector) representation as well; that one
% IS in closed form (Micro R, moffatt2007estimation; micro_R / micro_IR in the code) and is
% ruled out by cost, not by absence of a closed form. Keeping both is also what makes this
% paragraph match the two-closure box below.

\subsection*{The observable is an interval average}

A recorded sample is the current averaged over the acquisition window rather than sampled at an
instant. The current is low-pass filtered in analogue form before it reaches the converter, so
each stored value carries a weighted average of the recent current and not a reading of it at a
point. Writing the sampling period as $\Delta$ and taking the weighting to be uniform across it,
the quantity a likelihood must describe is the windowed average
\begin{equation}
  \bar{y}_t \;=\; \frac{1}{\Delta}\int_{t-\Delta}^{t} y(\tau)\,\mathrm{d}\tau \;+\; \eta_t,
  \label{eq:obs-avg}
\end{equation}
where $y(\tau)$ is the instantaneous population current and $\eta_t$ is the instrumental
noise accumulated over the window. This interval average is the physical observable. When
$\Delta$ is much shorter than the fastest relaxation of $\mathbf{Q}$ the average sits close
to the instantaneous current and the distinction is negligible; when $\Delta$ is comparable
to or longer than that relaxation the two diverge, and a likelihood built on instantaneous
sampling describes the wrong random variable.

Uniform weighting is a first approximation to what the electronics do. The effective measurement
kernel is the impulse response of the amplifier's analogue Bessel low-pass, whose tails reach into
the neighbouring intervals; the uniform window instead keeps the whole kernel inside one
acquisition interval, which is what lets a likelihood conditioned on the two ends account for the
averaging. A Bessel kernel also correlates successive recorded values through the signal and colours the noise that
Eq.~\ref{eq:obs-avg} takes to be white. The simulator and every member share the uniform window
(Methods), so data and likelihood describe the same observable and the window does not enter the
measurement here; the closures below are the approximation the diagnostic sees. The Discussion
prices the window, and carrying a general kernel through the same construction is left to later
work.
% CORRECTED 2026-08-04. The clause read "the closures below are the ONLY approximation under test",
% which is a rule-4 universal and is contradicted by 05_discussion.tex, where the uniform window
% against a Bessel kernel is conceded as a second approximation adopted knowingly. What keeps that
% one out of the measurement is that it is shared, which is what the sentence now says.
% FIXED 2026-08-06. The comment block above used to run on past its last line: "Carrying a general"
% ended the commented line and the next two lines carried no %, so they typeset as body text and
% the paragraph closed with the fragment "member share it. kernel, with the coloured noise ...".
% The sentence is restored whole as body text, which is what it reads as.
% COMPRESSED 2026-08-09 (length pass). Gone from this paragraph: the reason a Bessel filter is the
% one chosen (flat group delay, small step overshoot), which no later sentence uses; and the
% duplicate concession "The uniform window is itself an approximation to the impulse response of
% the anti-aliasing filter", which restated the paragraph's own opening sentence. All three parts
% of the caveat are still here and in this order: the uniform window approximates the amplifier's
% kernel; simulator and every member share it, so it is out of the measurement; a general kernel,
% with its coloured noise, is later work. The "last two sentences of the paragraph above" that
% [BESSEL-BOUND] below offers to replace are now the one beginning "The simulator and every member"
% and the one beginning "The Discussion prices".
%
% =========================================================================================
% [BESSEL-BOUND] OPEN 2026-08-06. THIS IS THE CANONICAL NOTE; eight other sites carry pointers
% to it and nothing else (grep -rn BESSEL-BOUND sections/). Read it here.
% =========================================================================================
%
% THE PROBLEM WITH THE PARAGRAPH ABOVE is not that it is wrong. It is that the manuscript
% concedes the Bessel kernel three times, here, at 06_methods.tex:73-77 and at
% 05_discussion.tex:38, and all three times WITHOUT A BOUND. An unbounded concession reads as
% "size unknown", which is a weaker position than the algebra supports and is the opening for the
% objection this scoping exists to answer: that the method does not apply to real recordings.
% The measured position is stronger and is available at no experimental cost.
%
% WHAT THE ALGEBRA GIVES. Computed, not asserted, in
% papers/1_method/decisions/recompute/bessel_bound.py (numpy + scipy, no data files, ~30 s).
% The real acquisition kernel is boxcar(Delta) composed with Bessel(f_c), and the likelihood
% implements the boxcar exactly, so the unmodelled part is governed by ONE dimensionless group,
% f_c*Delta. Its effect on the two objects the likelihood uses, the variance of a recorded value
% and the covariance between consecutive ones, for any component flat on the scale of Delta:
%
%   f_c*Delta |  0.2     0.5     1.0     2.0     5.0    10.0    20.0
%   deficit   | 61.5%   29.6%   14.7%    7.4%    2.9%    1.5%    0.74%
%   rho_1     | 64.1%   21.1%    8.6%    4.0%    1.5%    0.75%   0.37%
%   CI error  | 61%     19%      8.3%    3.9%    1.5%    0.74%   0.37%
%
% Four properties of that table matter more than the numbers in it.
%   (i) FIRST ORDER, AND NOTHING BEYOND LAG ONE. From f_c*Delta = 1 upward the deficit is
%       0.147/(f_c*Delta) and rho_1 is 0.075/(f_c*Delta), while rho_2 and every longer lag are
%       zero to four decimals. The residue is a nearest-neighbour effect of known size.
%  (ii) THE TOTAL IS CONSERVED. Var + 2*sum_k Cov_k = S/Delta to six decimals at every setting,
%       forced by the filter having unit gain at DC. The filter moves a fraction
%       0.147/(f_c*Delta) of each sample's variance into its two neighbours and destroys none of
%       it, so a likelihood that allows no such covariance sees a misallocation of variance
%       between the diagonal and its neighbours, of that size.
% (iii) POLE COUNT DOES NOT ENTER. Four poles, eight poles and a Gaussian of the same -3 dB
%       frequency agree to the third decimal from f_c*Delta = 0.5 upward, so the width can be
%       quoted as the Gaussian equivalent, sigma_f = sqrt(ln2)/(2 pi f_c) = 0.1325/f_c,
%       equivalently a 10-90 rise time of 0.34/f_c.
%  (iv) THE KINETICS ARE THE PROTECTED PART. The deficit is first order in 1/(f_c*Delta) for the
%       instrumental noise and SECOND order for a component carrying its own correlation time,
%       because unit DC gain preserves the integral of the autocovariance and the block-average
%       variance of a short-correlated component depends on that integral alone to leading
%       order. At f_c*Delta = 5 the white term loses 2.94% and a gating term with tau_c = Delta
%       loses 0.11%, a factor of 27. The residue therefore lands in the instrumental-noise
%       partition, which is exactly where the paragraph at 05_discussion.tex:35 already argues
%       an unmeasured departure lives. The two concessions are one concession.
%
% WHAT DOES NOT SCALE WITH THE GROUP, and it must not be buried: the interval holding the
% concentration jump. Group delay plus rise displace and smear the transient by about 0.69/f_c at
% four poles, a fixed time rather than a fraction of Delta, so that one interval carries an error
% the expansion says nothing about. The honest recipe discards or flags the first interval after
% a step.
%
% ONE THING THIS FAMILY GETS FREE that a point-sampling method does not. Because the observable
% IS the integral, block-averaging a record down to a coarser Delta produces exactly the
% observable of Eq. obs-avg at any decimation factor, with no aliasing question to answer. After
% decimation the only unmodelled piece left is the residual Bessel kernel, which is the table.
% This is what makes "average the record first" a scoping statement rather than a workaround.
%
% DRAFT BODY SENTENCES, replacing the last two sentences of the paragraph above if the bound is
% to be printed. About 90 words, no figure, no run:
%   "How much of the kernel the uniform window misses is set by one dimensionless group, the
%   acquisition window measured against the filter's response time. Because the filter has unit
%   gain at zero frequency it moves variance between a recorded value and its neighbours rather
%   than removing any, and the fraction it moves falls as the reciprocal of that group, about a
%   seventh where the window equals the filter's inverse cutoff and three per cent at five times
%   it, with nothing beyond adjacent values. A record digitised far above the cutoff is therefore
%   block-averaged to that scale before analysis, which yields exactly the observable of
%   Eq.~\ref{eq:obs-avg}. What averaging does not repair is the interval holding the
%   concentration jump, which the filter's group delay displaces by a fixed time."
%
% [OPEN, Luciano's call] Print the numbers here, or hold them for a referee reply. Either way the
% bound belongs in ONE place and Theory is it, because this is where the premise is made; the
% other sites should point rather than restate. Also open: whether to cite the Gaussian-filter
% rise-time convention (Colquhoun and Sigworth, Single-Channel Recording) for sigma_f, which
% would make (iii) a reader's own arithmetic instead of ours. That PDF is not in
% docs/bibliography/ yet, so nothing is cited from it here.

\subsection*{Two Gaussian approximations}
% CORRECTED 2026-08-06. The opening read "Every member of the family makes the first; it makes the
% second only if it predicts the interval average at all", and the interval-signal box read that a
% member predicting an instant "makes no interval-signal closure at all". Both sentences read as
% denying that the instantaneous members Gaussianize the observation, which is false: they model the
% recorded current as Gaussian just as the interval members do. What is true is narrower and is now
% what is stated. The instantaneous current is gamma^T applied to the occupancies at one time, a
% LINEAR functional, so the occupancy Gaussian pushes forward to it exactly and no second moment
% matching is performed; the interval average is a functional of the PATH inside the window and is
% not a function of the occupancies at any single time, so it needs a closure of its own. The
% instantaneous member's error at the signal level is therefore misspecification (a snapshot in place
% of the average of Eq. obs-avg, discarding the within-interval evolution), which is what
% 02_theory.tex:61 already calls describing the wrong random variable, and it compounds because
% averaging is the operation that makes the law more Gaussian. Line 34 was adjusted in step ("which
% object the second Gaussian describes").
% COMPRESSED 2026-08-09, length pass. What left this subsection: the etymology of "closure" and the
% fixed-dimension argument (both restated in the box's occupancy entry), the two "What is closed"
% entries of the box, the lumpability and population-convolution excursus, and the tilted-semigroup
% MGF sentence, which travels with Appendix 1. No claim, caveat or concession was dropped; the
% instantaneous-member point the note above protects is now carried by the opening paragraph alone,
% which is the only place it is stated, so it must not be cut from there.

To close a distribution is to replace it at that step by the Gaussian with the same mean and
covariance, discarding every higher moment. Two closures are available, on two objects. Every member of the family makes the first, and every member ends up with a Gaussian for the
recorded sample; what separates them is whether that Gaussian costs a second approximation. A member predicting the current at an instant does not need
it: that current is $\bm{\gamma}$ contracted against the occupancies at one time, and the linear
image of a Gaussian is Gaussian. A member predicting the interval average does, since the average
is a functional of the path inside the window rather than of the occupancies at any one time. What
an instantaneous member drops is the evolution of the channels inside the window, one snapshot
standing in for the average the recording holds, which misspecifies the random variable rather than
closing it, and the substitution runs in the unhelpful direction, since averaging over the window is
itself what brings the law closer to Gaussian.
Matching moments at each step does not keep them true across a recording: the update a closure
licenses is itself approximate, so the carried pair drifts from the exact one.

\begin{tcolorbox}[title={The two Gaussian approximations}, colback=gray!5, colframe=gray!55]
\textbf{Occupancy closure.}
\emph{Exact object:} one probability per occupancy vector, a count that grows with
$N_{\mathrm{ch}}$.
\emph{Replacement:} a Gaussian in the per-channel pair $(\bm{\mu},\bm{\Sigma})$, the same size at
every channel count; because the occupancies sum to one, $\bm{\Sigma}$ is singular and the Gaussian
lives on the simplex.
\emph{Validity:} two conditions, not one. The channel count must be high enough for the counts to
behave as a continuum, which is the central limit theorem and fails as $N_{\mathrm{ch}}$ falls;
and the recorded sample must not resolve the step in predicted current made by moving one channel
between states, which fails as the instrumental noise falls. The second is why the boundary of the
failing region in the Results runs across the plane of channel count against noise rather than
down a line of constant $N_{\mathrm{ch}}$.

\smallskip
\textbf{Interval-signal closure.}
\emph{Exact object:} the law of the conductance averaged over one acquisition window,
telegraph-like and with no closed form (below).
\emph{Replacement:} a Gaussian carrying its first two moments.
\emph{Validity:} the window must span many transitions, or the instrumental noise must be large
enough to make the departure immaterial. The second half is exact rather than heuristic: noise
independent of the gating current adds to the variance of the recorded sample and leaves every
cumulant above the second unchanged, so the standardized departures shrink as the noise grows.

\smallskip
The two are controlled by different quantities, the occupancy closure by the channel count and
the interval-signal closure by the window measured against the gating timescale, while
instrumental noise relaxes both by making the recorded sample more Gaussian. That makes them separable by design, and the Results sweep a
plane of channel count against noise at one window before moving the window along an axis of its
own.
\end{tcolorbox}

\emph{Telegraph-like} names this object: a single channel switching among its conductance levels
traces a piecewise-constant path, the random telegraph signal when $K=2$, and the window delivers
its integral over one acquisition interval,
\begin{equation}
  A_{0\to\Delta}=\int_0^{\Delta}\gamma\bigl(X_s\bigr)\,\mathrm{d}s,
  \label{eq:integrated-conductance}
\end{equation}
where $X_s$ is the state of the channel at time $s$. Its law is mixed: an atom at
$\Delta\gamma_i$ carrying the probability of never leaving state $i$, and a continuous part from
the paths that switched. Its density is a series over the number of transitions in the window,
exact at every $K$, closed only at $K=2$, and proliferating with $K$. Having few distinct
conductances does not rescue it, which is worth saying because channel models are exactly of that
kind: every shut state carries $\gamma=0$ and the open states often share one $\gamma$, so
$A_{0\to\Delta}$ is the open conductance times the time spent in the open aggregate however large
$K$ is, and that aggregate is Markov only if $\mathbf{Q}$ is lumpable, which for the schemes that
motivate this work it is not \citep{colquhoun1981stochastic}. The population multiplies the
difficulty rather than adding to it: the laws of independent channels convolve, and a recorded
sample constrains only the sum, so the posterior over channels does not
factorize~\citep{moffatt2007estimation} and an exact recursive filter has to carry a distribution
over joint occupancy vectors.
% TODO citation: the two-state occupation-time density is classical (candidates to check and
% add to biblio.bib: Takacs 1957, Acta Math Acad Sci Hungar; Pedler 1971, J Appl Prob 8:381).
% Neither PDF is in docs/bibliography/ yet. The general-K series is the accumulated-reward
% problem of performability analysis (Donatiello and Iyer 1987; de Souza e Silva and Gail
% 1986; Sericola 1990s). Cut 2026-08-06 with the uniformization/B-spline paragraph.
% OPEN 2026-08-06. A clause here pointed at the Discussion for the other consequence of the same
% aggregation, that distinct generators can produce identical current statistics. It was removed
% because 05_discussion.tex carries no such passage: grep for aggregat/lumpab/identifiab across
% sections/ returns only this paragraph. The open item is already logged at 00_abstract.tex:203
% ("worth a Discussion sentence with the identifiability literature"); the sourced treatment is in
% docs/bibliography/identifiability/ and the framing in _program/research_program.md (Bridge 2).
% If that sentence gets written, restore the cross-reference here.
% CONTEXT 2026-08-09: the aggregation passage that clause sat inside (few distinct conductances,
% the lumpability caveat with colquhoun1981stochastic, the population convolution and the
% non-factorizing posterior of moffatt2007estimation) went in the length pass, so the note has no
% host sentence now; the open item stands. What the asymmetry sentence below needs from that
% passage, the proliferation in N_ch, is folded into the density sentence above.
% colquhoun1981stochastic survives at 01_introduction.tex:54 and moffatt2007estimation is cited
% elsewhere in this section, so no bibliography entry is orphaned.
The closure is indifferent to all of it: it discards every cumulant above the second, and the
moments it keeps are closed-form single-channel objects, independent of $N_{\mathrm{ch}}$ and of
the topology, two derivatives of one $K\times K$ matrix exponential
(Appendix~\ref{app:derivation}). That asymmetry, an exact law that degrades with $K$ and with
$N_{\mathrm{ch}}$ against moments that do neither, is what justifies the closure beyond the
two-state case run here.

The higher moments these closures discard do not vanish. They reappear in the data as temporal
correlation the likelihood does not predict, which is what the correlation term of the
information-distortion decomposition measures (Diagnostics). Three regimes follow, named here so
the maps in the Results have vocabulary: the \emph{microscopic} regime, few channels, where the
occupancy closure fails and only the microscopic description is exact; the \emph{telegraphic}
regime, windows far shorter than the relaxation time, where the interval-signal closure fails;
and the \emph{Gaussian} regime, where both closures hold. Any Gaussian likelihood of this kind
must fail somewhere. The aim of this paper is to show that the failures are the predicted
degradation of these two closures rather than defects of a particular algorithm.

% MOVED 2026-08-06 to 04_results.tex, into the paragraph that reads the five regions in order.
% The paragraph that stood here mapped the three regimes onto the regions of the usage map. Half of
% it was already said twice (the negative slope, in the occupancy-closure box above at "why the
% boundary of the failing region in the Results runs across the plane", and again in the Results
% prose), and the half that was not could not be read here: it named four regions by their contents
% before the reader has seen the figure. The three regimes are still named in the paragraph above,
% which is what the Results vocabulary needs from this section.

\subsection*{Least squares, and the two axes of the family}
% RETITLED 2026-08-06, twice in one pass. "The ladder below it" named the window axis alone and
% left the recursion axis out of the heading of the subsection that introduces it. A first pass
% called it a grid; the heading now names the axes rather than the shape, because the shape is a
% partial order and not a chain, and "grid" oversells a $4\times2$ of which two cells are empty.
% The cost ladder survives as what it always was, an ordering of computational cost, and stays in
% the Results heading and the abstract.
% RETITLED AGAIN 2026-08-06 (Luciano). "The root question, and the two axes below it" was the only
% heading in this section that named the rhetoric of the argument ("the root question") and a
% position on the page ("below it") instead of the object under discussion; the other four name
% their object. After the previous retitle "below" had also stopped describing anything: the two
% axes do not hang off the root question, they are the two knobs the root question leaves outside.
% The heading now names the two things the subsection does, in the order it does them.

Before any of the choices above comes a prior question: does the method model the gating
fluctuations at all? Classical nonlinear least squares answers no. It fits
the deterministic mean current and assigns one constant variance to every residual, so it
carries neither an occupancy distribution nor a conductance conditioned on anything, and
the temporal structure the two closures approximate is absent from it. Least squares is
therefore not a member of the family. It is the bottom rung of the ladder of computational
cost this paper walks, and it is the anchor against which every member above it is
measured, because it is the method most readers are using.

The methods differ along two axes. One axis is the window: how much of the acquisition
interval the predicted current is conditioned on. The other is recursion: whether the
occupancy distribution is conditioned on the data as the recording proceeds.

The window axis is not the family's property alone. What the amplifier delivers is an
average over the acquisition interval, so any method that predicts a recorded value has to
decide whether to predict the current at an instant or the current averaged over the
window, and that decision is available at every level of the gating description, least
squares included. Least squares therefore has two arms here. They carry one constant
variance and are otherwise identical, so the pair isolates the acquisition average with
nothing else moving, which no pair inside the family does.

Along the recursion axis, a recursive (prequential) member factorizes the likelihood as
$\prod_t q_t(\bar{y}_t \mid \bar{y}_{1:t-1})$, so the occupancy distribution carried into
each window is the posterior left by the previous windows and the hidden state is
conditioned on the data. A non-recursive (open-loop) member factorizes as
$\prod_t m_t(\bar{y}_t)$, with each window's predictive propagated from the initial
condition and never conditioned on the observed trace.

The members are named from the two axes and the names are read off them. The suffix is the
recursion axis: \texttt{NR} for a non-recursive member, \texttt{R} for a recursive one. The
% CORRECTED 2026-08-06 (Luciano). The suffix was given as "N or R", which is how
% papers/_program/nomenclature.md:98 states it and which does not parse: NR ends in R, so by that
% reading the one member the letter is meant to mark as non-recursive comes out recursive. The
% split that works on all six names is NR against R, since NR and INR end in the first and R, MR,
% VR and IR in the second. nomenclature.md:98 was aligned the same day and carries the same note,
% so the two are in step; keep them that way.
prefix is the window axis, and it says how much of the interval the conductance is
conditioned on: no prefix for the current at an instant, \texttt{M} for the interval mean
given the state at the window's start, \texttt{I} for the interval mean given the states at
both of its ends. Five members follow rather than six, because without an update the state
at the end of a window is never observed and the two interval prefixes describe the same
method on the non-recursive side: \texttt{NR} and \texttt{R} predict an instant, \texttt{INR}
is the non-recursive interval member, and \texttt{MR} and \texttt{IR} are where the start
state and the boundary separate. The \texttt{M} is the mean conductance and not an
abbreviation of macroscopic; every member of the family is macroscopic. A third choice, open
only to the members that stop at the start state, is the form given to the interval variance,
and it accounts for the one letter that names neither axis: \texttt{VR} is \texttt{MR} with
the residual interval variance in place of the total one, the two differing as set out at
Eq.~\ref{eq:mr-vr} below. Least squares carries no gating description at all and sits off
this lattice, and we write \texttt{LSE} for the fit to the instantaneous deterministic mean
and \texttt{ILSE} for the fit to the same mean averaged over the interval, where the
\texttt{I} records the window setting and nothing else. Eight in all: five on the lattice,
\texttt{VR} beside \texttt{MR}, and the two least-squares arms. The window decision is what
separates \texttt{LSE} from \texttt{ILSE}, \texttt{NR} from \texttt{INR}, and \texttt{R} from
\texttt{MR} and \texttt{IR}.
% ADDED 2026-08-06 (Luciano), and it closes the section's largest gap. Before this paragraph the
% eight codes were USED before being defined anywhere the reader had reached: NR, INR, R and IR
% first appeared in the sentence that coined LSE and ILSE, as pairing partners with no content, and
% MR and VR appeared 200 lines later inside the av=1 against av=2 algebra. The only definitions in
% the section were Table 1 at its end, with a forward pointer, so six empty codes had to be carried
% through "The boundary state" and "The interval update", which are exactly where MR/VR/IR is
% decided. The other two definitions arrive later still: the Figure 1 caption glosses six of the
% eight (not MR, not VR) and the flag table is in Methods, which in eLife order follows Results.
% What is stated here is the NAMING RULE rather than eight glosses, so the codes become legible
% instead of memorised; the rule is papers/_program/nomenclature.md:98-121 (suffix = occupancy
% axis, prefix = conductance axis, V = variance form) and had never been written into the
% manuscript. Three further repairs ride on it: (i) the old pairing clause read "R pairs with IR",
% which skips MR, the member the naming rule puts between them, so the pairing is now stated over
% all three; (ii) "Eight members follow from the two axes" was arithmetic the axes do not support,
% since VR comes off a third axis and the two least-squares arms are off the lattice, which Methods
% says (06_methods.tex:176) and Theory did not; (iii) a reader who has met MacroIR in the abstract
% can read the M of MR as "Macro", which nothing disarmed.
% PLACED after both axis paragraphs rather than between them, so the rule is stated once both
% axes it reads from have been given, and it lands immediately before the pointers to Figure 1 and
% Table 1, which are what the reader checks it against.

These are the two closures seen from the side of the method rather than the side of the
approximation. What the occupancy closure is applied to is set by the recursion axis, and
what the interval-signal closure is applied to is set by the window axis. One conditioning
choice governs both, and this is worth stating before the algebra rather than after it: the
states the interval is conditioned on are the states the interval-averaged conductance is
computed from \emph{and} the states the Bayesian update is taken on, so moving along the
window axis moves the conductance and the gain together. Both members of that statement are
set out at the end of this section, once the equations they differ in are on the page:
Table~\ref{tab:members} places all eight against those equations, where the axes can be read
off as settings rather than as a taxonomy, and Figure~\ref{fig:ladder} draws six of them as
one turn of the filter cycle each.
% RESTRUCTURED 2026-08-06 (Luciano). What stood here was a two-dimensional table introducing
% the eight members before the reader has met mu, sigma^2_pred or the gain, which is the job
% Figure 1 already does better and in pictures. Two axes are two sentences; what needs a table is
% the anatomy of the predictive Gaussian member by member, and that table now sits after the
% equations it indexes. Also dropped with it: the word "grid" for the design. Two knobs, each
% monotone in cost, give a partial order and not a chain, which is the same fact the Results
% report when the calibration comes out non-monotone.

Two things follow from the window axis. First, \texttt{IR} occupies the top rung below the
intractable exact likelihood, so its standing in the family is structural and set before
any measurement. Second, the lower rungs are established likelihoods rather than
constructed foils: the non-recursive instantaneous member is the independent-interval
likelihood of \citet{milescu2005maximum}, and the recursive instantaneous member is the
macroscopic recursive filter of \citet{moffatt2007estimation} and the generalized filter
of \citet{munch2022bayesian}. The lattice is one coordinate system that contains these
methods as special cases.

\subsection*{The boundary state}

A \emph{boundary state} is the pair $(i\to j)$ of a channel's states at the two ends of
an acquisition interval. Chaining intervals, the end state of one interval is the start
state of the next, so there is a single state variable per junction: the filter conditions
on the states at the interval boundaries and marginalizes the trajectory between them
analytically. This is the same device as static condensation, or the spatial Markov
property, applied on the time axis. It is not a transition state in the mechanistic sense,
and the word \emph{transition} is not used for it here. Because the joint dynamics of the two
endpoints are generated by the Kronecker sum
$\mathbf{Q}\oplus\mathbf{Q}=\mathbf{Q}\otimes\mathbf{I}+\mathbf{I}\otimes\mathbf{Q}$, whose
propagator factorizes as $e^{(\mathbf{Q}\oplus\mathbf{Q})t}=e^{\mathbf{Q}t}\otimes
e^{\mathbf{Q}t}$, the moments over the $K^2$ boundary pairs collapse to bilinear forms in
the $K$-dimensional state space and no explicit $K^2$ object is formed (Methods).
\subsection*{The interval update}

The interval likelihood is the instantaneous one applied to the boundary state. The recursive
member of \citet{moffatt2007estimation} carries a Gaussian over the channel occupancies and
conditions it on each recorded sample: in the per-channel normalization used throughout, with
$\bm{\gamma}$ the per-state conductance and $\epsilon^2$ the instrumental variance,
\begin{align}
  y^{\mathrm{pred}} &= N_{\mathrm{ch}}\,\bm{\mu}^\top\bm{\gamma},
  \qquad
  \mathbf{g} = \bm{\Sigma}\bm{\gamma},
  \qquad
  \sigma^2 = \epsilon^2 + N_{\mathrm{ch}}\,\bm{\gamma}^\top\bm{\Sigma}\bm{\gamma},
  \nonumber\\
  \bm{\mu}^{\mathrm{post}} &= \bm{\mu} + \frac{\mathbf{g}}{\sigma^2}\,\delta,
  \qquad
  \bm{\Sigma}^{\mathrm{post}} = \bm{\Sigma} - \frac{N_{\mathrm{ch}}\,\mathbf{g}\mathbf{g}^\top}{\sigma^2},
  \qquad
  \delta = y^{\mathrm{obs}} - y^{\mathrm{pred}}.
  \label{eq:macror}
\end{align}
This is the Gaussian conditioning of the pair (state, observation) on the value recorded. The
gain $\mathbf{g}$ is the covariance between the state a channel is in and the current it
carries, and it is the whole of what brings the observation to the state: the predicted
variance divides it, and nothing else in the update touches the data. It is also the object
the interval version has to rebuild, since the two contractions of Eq.~\ref{eq:macror} against
$\bm{\Sigma}$ are $\mathbf{g}$ itself and $\bm{\gamma}^\top\mathbf{g}$, the same product taken once
and twice.
% SLOT SYMBOLS, settled 2026-08-06 (Luciano), after two rounds of getting it wrong. The first draft
% of this equation wrote it for "a state space of any size" with neutral symbols, c for the
% conductance and k for the gain. Both were mistakes. k collides twice over, with the state count K
% one case away and with the channel index in N = sum_k X_k in the next paragraph. And c was
% unnecessary: this equation is not a template, it is MacroR, which is published with gamma
% (elife-template.tex:360-384), and the tilde already carries gamma inside it, so a neutral letter
% put two symbols on one slot within a page for nothing. The move the three steps make is a
% SUBSTITUTION into these equations, not an instantiation of a template, which is also how Luciano
% states it. The gain gets no symbol at all here, exactly as in the published form: it is written
% Sigma*gamma inline. FINAL FORM (Luciano's call): Sigma*gamma does get a dedicated symbol, and it
% is g, the name the object already had at Eq. mu-post and in the derivation, introduced here where
% it is first built instead of arriving unexplained thirty lines later. One symbol, one meaning
% across the three members: the covariance between the state that will be carried forward and the
% recorded current. Its K^2 form on the pairs is written inline and never named, so g always denotes
% the K-vector the update consumes. This also cost a cleanup at 06_methods.tex:701, which used a
% bare undefined $g$ for a unitary conductance.
Every moment in this section belongs to a single channel. Writing $\mathbf{X}$ for the
indicator vector of the state one channel occupies, which is one on that state and zero
elsewhere, $\bm{\mu}=\mathbb{E}[\mathbf{X}]$ is its probability vector and
$\bm{\Sigma}=\mathrm{Cov}(\mathbf{X})$ its covariance, which is
$\mathrm{diag}(\bm{\mu})-\bm{\mu}\bm{\mu}^\top$ for a channel nothing has been observed about
and vanishes when its state is known. Both are objects of one channel, and neither is a
population quantity divided by anything. Channels are independent of each other, so the counts
$\mathbf{N}=\sum_c \mathbf{X}^{(c)}$ over the ensemble have
$\mathbb{E}[\mathbf{N}]=N_{\mathrm{ch}}\bm{\mu}$ and
$\mathrm{Cov}(\mathbf{N})=N_{\mathrm{ch}}\bm{\Sigma}$, and that is where every factor of
$N_{\mathrm{ch}}$ in Eq.~\ref{eq:macror} comes from: the recorded current is a population
quantity and the state it is informative about is carried per channel.

\paragraph{What Eq.~\ref{eq:macror-w} gets exactly right, and what it gives up.}
The three moments the update consumes are exact functions of $(\bm{\mu},\bm{\Sigma})$. In
particular the per-state variance enters $\sigma^2$ and touches neither the predicted mean nor
$\mathbf{g}$: given the counts, the extra fluctuation has mean zero, so
$\mathrm{Cov}(\mathbf{N},\bar y)=\mathrm{Cov}(\mathbf{N},\mathbb{E}[\bar y\mid\mathbf{N}])$ and
$\mathbf{w}$ drops out of it. That is what makes charging the leftover variance to the noise term
legitimate rather than convenient, and it is used at Eq.~\ref{eq:pred-var} below.

What is given up is not exactness but \emph{sufficiency}. Because
$\mathrm{Var}[\bar y\mid\mathbf{N}]$ depends on $\mathbf{N}$, the prior over the state enters the
variance of the likelihood and not only its mean, and a normal prior is conjugate for the mean of a
normal likelihood but not for its variance. The pair $(\mathbf{N},\bar y)$ is therefore not jointly
Gaussian, and the update of Eq.~\ref{eq:macror} is not the exact Bayesian update: it conditions a
Gaussian that carries the exact first two moments, treating a heteroscedastic emission as if it
were homoscedastic. \citet{munch2022bayesian} meet the same obstruction and resolve it the same
way, noting that the state-dependent term makes the incremental likelihood intractable because
``the conjugacy property is lost'', applying the total-variance decomposition to recover mean and
covariance, and stating the cost in the terms used here: the covariance equations ``are still exact
for the more general noise model'', while ``mean and covariance just do not form a set of sufficient
statistics anymore''. Two groups reached the same device from different motivations, which is worth
more to a reader than either group's assurance.

% NEW PROSE, 2026-08-09, Phase 2 of papers/1_method/08_length_plan.md, and the only new writing in
% the pass. It replaces about 4,800 words of derivation that moved to 08_appendix_derivation.tex.
% WHAT IT IS FOR: a reader who will not open the appendix should still know how the members were
% built; a reader who will should know what is there. It states the construction, not its merits
% (rule 5), and it says which steps are exact and which is the approximation rather than claiming
% exactness for the whole (the appendix proves the split).
% IT ALSO CARRIES A DEFINITION THE BODY WOULD OTHERWISE LOSE. The tilted semigroup and
% $\bm{\Gamma}=\mathrm{diag}(\bm{\gamma})$ were defined once, inside the telegraph paragraph that
% the closures compression cuts. Two lines of the appendix say "met above" and "as above" and both
% resolve here. Do not delete the definition without repairing 08_appendix_derivation.tex:46 and :56.
% IT MUST NOT DUPLICATE the sufficiency paragraph above it, which already says the update conditions
% a Gaussian carrying the exact first two moments and is not the exact Bayesian update.
What stops Eq.~\ref{eq:macror} from being used on the interval average is that the average is not a
function of the occupancy at any one time, so no conductance vector on the $K$ states delivers its
mean. On the $K^2$ pairs of states a channel occupies at the two ends of the window one does, and
the construction of every interval member is that substitution carried out in three steps. Run the
update on the pairs, with the pair-conditioned mean interval conductance $\bar\gamma_{i\to j}$ in
place of $\bm{\gamma}$. Charge what the conductance still varies by once the pair is fixed, its
variance $\bar v_{i\to j}$, to the noise term of Eq.~\ref{eq:macror-w}, which is legitimate because
that fluctuation has mean zero given the counts and so enters neither the predicted mean nor the
gain. Then marginalize the start index, which the next window has no use for.

The two moments this asks for are cheap at any $K$. Both are derivatives at $\alpha=0$ of the tilted
semigroup $e^{(\mathbf{Q}+\alpha\bm{\Gamma})\Delta}$, with
$\bm{\Gamma}=\mathrm{diag}(\bm{\gamma})$, so each costs one matrix exponential and no object larger
than $K\times K$ is formed in building them. Of the three steps the first two are exact, and the
Gaussian replacement is where the approximation enters. Appendix~\ref{app:derivation} carries the
derivation, the boundary covariance the pairs induce, the spectral form the implementation
evaluates, and what separates the members that stop at the window's start from the one that
conditions on both of its ends.

In a recursive member the posterior pair $(\bm{\mu}^{\mathrm{post}}, \bm{\Sigma}^{\mathrm{post}})$
becomes the prior for the next window; in a non-recursive member it is not carried forward,
and each window's predictive is propagated from the initial condition alone. Reusing the
posterior as the next prior is itself a closure: it is exact only while the end marginal is
a sufficient summary of the boundary posterior, which the Gaussian projection does not
guarantee. For this reason the recursive members' Gaussian information is a model-based
quantity and is not claimed to be the exact information of the process; how far it departs
from the truth is what the Diagnostics measure.

\begin{table}[htbp]
\begin{fullwidth}
\centering
\begin{tabularx}{\linewidth}{@{}l
    >{\hsize=0.80\hsize\raggedright\arraybackslash}X
    >{\hsize=0.95\hsize\raggedright\arraybackslash}X
    >{\hsize=1.25\hsize\raggedright\arraybackslash}X@{}}
\toprule
Member & The states are conditioned on & Predicted mean (Eq.~\ref{eq:ypred}) & Predicted variance beyond the instrumental term $\epsilon^2_{0\to\Delta}$ (Eq.~\ref{eq:pred-var}) \\
\midrule
\multicolumn{4}{@{}l}{\emph{Open loop}: no update; every window is propagated from the initial condition (Eqs.~\ref{eq:mu-prop},~\ref{eq:sig-prop} alone)} \\
\addlinespace[2pt]
\texttt{LSE} & nothing: there are no states & the deterministic mean current at the middle of the window & none: one constant over the whole record, fitted as the residual sum of squares \\
\texttt{ILSE} & nothing: there are no states & the same mean current averaged over the window & the same constant \\
\texttt{NR} & an instant, the middle of the window & the current at that instant & the spread of the instantaneous conductance across the states the channels occupy \\
\texttt{INR} & the state at the start of the window & the interval-mean current given that state & the spread of that interval mean across the occupied states, and the variance of one channel's interval mean given its start state \\
\midrule
\multicolumn{4}{@{}l}{\emph{Recursive}: the posterior of one window is the prior of the next (Eqs.~\ref{eq:mu-post},~\ref{eq:sig-post})} \\
\addlinespace[2pt]
\texttt{R} & an instant, the middle of the window & the current at that instant & the spread of the instantaneous conductance across the states the channels occupy \\
\texttt{MR} & the state at the start of the window & the interval-mean current given that state & the spread of that interval mean across the occupied states, and the variance of one channel's interval mean given its start state \\
\texttt{VR} & the start, and both ends for the interval variance alone & the interval-mean current given the start state & the spread of that interval mean across the occupied states, and the variance of one channel's interval mean given both its endpoints \\
\texttt{IR} & the states at both ends of the window & the interval-mean current, the unobserved end state marginalized & the spread of the interval mean across the occupied pairs of endpoint states, and the variance of one channel's interval mean given both its endpoints \\
\bottomrule
\end{tabularx}
\end{fullwidth}
\caption{\textbf{The eight members as settings of the equations above.} One conditioning
choice runs through each row: the states an interval is conditioned on are the states its
conductance is computed from and the states its update is taken on, so the second column
governs the gain of Eq.~\ref{eq:mu-post} as well as the two columns beside it. Reading down,
the second column is the window axis and the two blocks are the recursion axis. Three
consequences are worth reading off. \texttt{LSE} and \texttt{ILSE} differ in the mean alone,
which isolates the acquisition average with nothing else moving; \texttt{NR} and \texttt{R},
and \texttt{INR} and \texttt{MR}, carry identical Gaussians and differ only in whether the
recording is fed back into the states; and \texttt{MR} and \texttt{IR}, taken from the same
state, assign the same variance by two decompositions that rearrange into each other, so
what separates them is the gain, which is then what makes the variances they report along a
recording differ after all, while \texttt{VR} drops the between-endpoint spread from the
third term without restoring it in the second and predicts less than either from any
state. Read against Eq.~\ref{eq:emission}, that last row is not a fourth reading of the
construction but a mismatch: naming the state fixes the per-state variance, and \texttt{VR}
fills the slot with the leftover of a state it does not otherwise use, which is why it is the
one member whose variance form is a choice rather than a consequence. In print, \texttt{IR} is
MacroIR and \texttt{INR} is MacroINR. All eight are implemented and measured: six of them are
drawn as one turn of the filter cycle in Figure~\ref{fig:ladder}, seven are run over the
design plane and the un-averaged least-squares arm at a single cell of it, and the
dispatched flags and data keys are in Table~\ref{tab:family}.} % src: papers/_program/nomenclature.md (the lattice as implemented; what the letters mean) ; legacy/qmodel.h:4531-4645, read 2026-08-06 for the three columns member by member: t_P is P_half at av=0 and P otherwise (mid-window evaluation state, qmodel.h:4489-4505), t_gmean_i is g at av=0 and gmean_i at av>0 INCLUDING av=2, gSg gains its boundary term only at av=2, and ms switches total against residual on (av==2 || variance_form==residual). NOTE for the record: the predicted MEAN of IR is built from gmean_i, the start-conditioned interval mean, exactly as MR's is; the boundary object enters the variance and the gain, not the mean. The row says "the unobserved end state marginalized" for that reason.
\label{tab:members}
\end{table}

% MOVED HERE FROM 04_results.tex 2026-08-06 (Luciano), with the subsection that framed it
% ("One filter step along the cost ladder") deleted rather than moved. The figure carries no
% measured claim by its own caption, it fixes the vocabulary, and Theory names it twice, so a
% display item that does no measuring was opening the Results and being cited two sections
% before the reader could see it. Division now: this figure is the mechanism (where inside one
% interval each choice acts), Table 1 is the specification (what each member puts into the mean,
% the variance and the gain), and Table 2 in Methods is the dictionary.
% ORDER, decided 2026-08-06 (Luciano) after a first placement up at the two-axes paragraph: the
% TABLE COMES FIRST and the figure follows it. The table cannot move up, because every equation
% it indexes (eq:ypred, eq:pred-var, eq:mu-post, eq:mr-vr) is inside "The interval update", so
% table-first means figure-last. That also repairs the argument for the earlier placement, which
% was that the figure is a map through the algebra: its rows ARE mu, sigma^2_pred, the gain and
% the score, the symbols the algebra defines, so the picture cannot precede them. Placing it up
% there also cost a full page of float in the middle of the two-axes argument and pushed Table 1
% onto page 15, a page into Diagnostics (measured on the 17:03 build).
\begin{figure}[htbp]
\begin{fullwidth}
\centering
\includegraphics[width=\linewidth]{Figure_1}
\end{fullwidth}
\caption{\textbf{One filter step along the cost ladder.} The current from a patch of
$N_{\mathrm{ch}}=20$ two-state channels, recorded as an interval average with white noise, is scored
by six members (columns). They form three pairs, one per level of the gating description: least
squares on the deterministic mean current (\texttt{LSE}), the open-loop gating likelihood
(\texttt{NR}) and the recursive filter (\texttt{R}), each beside its interval-averaged partner, named
with an I prefix (\texttt{ILSE}, \texttt{INR}, \texttt{IR}). Rows follow one turn of the cycle:
(A) the prior open probability from the Markov step, (B) the predicted current with its spread and
the innovation, (C) the Bayesian posterior, absent from the four open-loop columns, and (D) the
log-likelihood accumulated from zero. Row (B) is where the members of a pair part company: the
left-hand one predicts an instant, a point at the sample with an error bar, the right-hand one the
interval average, a segment across the interval with a band. In (A) and (C), \texttt{IR} alone
carries a pair of occupancies, one at each end of the interval, tied by a dashed arc: the arc is
their covariance and implies no path through the interval. The two-tone disc is the instant that
closes one interval and opens the next, and the Markov arrow leaves it because that marginal is all
that crosses the boundary. Colour encodes the filter phase (green the channel current, orange
predict, blue observe, purple update, black the score), arrows are operations, and a state carried in
from the neighbouring interval is faded. The two partial interval corrections, \texttt{MR} and
\texttt{VR}, sit off this pairing and are not drawn (Appendix~\ref{app:members}). No number is quoted from this 12 ms
recording.} % src: projects/eLife_2025/figures/paper_both/figure_1.Rmd (roster of six, 2026-08-05: LSE/ILSE/NR/INR/R/IR). NO NUMBER IS QUOTED FROM THIS FIGURE AND NONE SHOULD BE: the run is seeded with seed = 0, the sentinel for a random draw (ops/local/figure_1_plus_lse.macroir:75 and its own note at line 12), so the trace differs on every re-run and with it every quantity read off the realisation. The least-squares variance is SSE/n and moved from 0.67 to 12.5 between two runs of the same script for exactly that reason. The occupancy trajectory of row (A) is the deterministic Markov mean and does NOT move: verified on the rebuilt dumps, LSE and NR agree to seven decimals at every interval and so do ILSE and INR. That coincidence runs ACROSS the gating levels at one arm, not between the two members of a pair, which sit at different instants and cannot coincide; the caption no longer claims it, the version that did had it backwards. IR's boundary-pair marks (dashed tie, two-tone shared instant, where each operation arrow is anchored and why) are specified in figure_1_panels.R, the ir_* block. LENGTH: this caption is measured, not guessed. At 280 words the float spans 84 to 711 pt against a text block of 57 to 721, so it fits with ~19 pt of slack, under two lines. Adding 40 words overflows the page and the figure-supplement line lands on the footer, which is the failure the earlier 7.5 in note describes. Re-measure with a probe document on elife.cls before growing it. papers/_program/provenance.md#4 (20 channels, 250 substeps)
\label{fig:ladder}
% NOTE: the graphics file is still named Figure_S1.pdf on disk, from the retired independent
% supplementary numbering. It is written by figure_1.Rmd on the merged roster and is current;
% rename it to Figure_1_supplement_1.pdf in the deposit pass.
\figsupp[The same filter cycle over the first ten milliseconds of the recording.]{The same filter cycle over the first ten milliseconds of the recording, so that the concentration jump and the onset of gating are inside the window.}{\includegraphics[width=\linewidth]{Figure_S1}}
\end{figure}

\subsection*{The top of the ladder}

The exact likelihood conditions on the full trajectory inside each window and is
intractable, which is why the two closures above are needed at all. The same stochastic
model that makes the exact likelihood intractable can be simulated forward exactly, one
trajectory at a time, so the simulation realizes the object at the top of the ladder and
serves as the ground truth against which every rung below it is judged.
